\documentclass[11pt]{article}
\usepackage[margin=1.05in]{geometry}
\usepackage{amsmath,amssymb,amsthm,mathtools}
\usepackage{microtype}
\usepackage[hidelinks]{hyperref}

\newtheorem{theorem}{Theorem}
\newtheorem{lemma}[theorem]{Lemma}
\newtheorem{proposition}[theorem]{Proposition}
\newtheorem{remark}[theorem]{Remark}

\newcommand{\Q}{\mathbf Q}
\newcommand{\Z}{\mathbf Z}
\newcommand{\Qtwo}{\mathbf Q_2}
\newcommand{\vTwo}{v_2}
\newcommand{\Dn}{D}
\newcommand{\phiG}{\varphi}

\title{A modular--$2$-adic two-row approximation to Catalan's constant\\
with worthiness $0.9025266028\ldots$}
\author{}
\date{}

\begin{document}
\maketitle

\begin{abstract}
Let $G=\beta(2)$ be Catalan's constant and $\phiG=(1+\sqrt5)/2$.  We combine
Zudilin's Ap\'ery-like Catalan recurrence with a second, modular recurrence
coming from a weight-$3$ Eisenstein series with character $\chi_{-4}$.  The
key arithmetic input is that the two rational approximants converge to the
same $2$-adic Catalan period, but with different speeds.  Sampling the
Zudilin row at $5n$ and the modular row at $8n$ balances these two $2$-adic
speeds and produces a hidden determinant divisor of size $2^{40n+o(n)}$.
A correlated rank-two congruence lattice then gives rational approximants
$p_n/q_n$ satisfying
\[
 \liminf_{n\to\infty}
 \frac{-\log|G-p_n/q_n|}{\log|q_n|}
 \ge
 \frac{50\log\phiG}
 {10-2\log2+\frac{75}{2}\log\phiG}
 =0.902526602856971\ldots .
\]
The bound is below $1$, so it is not an irrationality proof.  We also explain
why the ratio $5:8$ is forced by the arithmetic and is optimal within this
particular two-row scalar architecture.
\end{abstract}

\section{The two Ap\'ery rows}

We write $D_N=\operatorname{lcm}(1,\dots,N)$.  Zudilin's Catalan sequences
$Q_m,P_m$ are characterized by
\[
 Q_0=1,\quad Q_1=\frac74,\qquad
 P_0=0,\quad P_1=\frac{13}{8},
\]
and the same second-order recurrence
\begin{align*}
 &(2m+1)^2(2m+2)^2(20m^2-8m+1)u_{m+1}\\
 &\quad=
 (3520m^6+5632m^5+2064m^4-384m^3-156m^2+16m+7)u_m\\
 &\qquad +(2m-1)^2(2m)^2(20m^2+32m+13)u_{m-1}.
\end{align*}
We use the standard asymptotics
\begin{equation}\label{eq:Zarch}
 \frac1m\log|Q_m|\to5\log\phiG,
 \qquad
 \frac1m\log|Q_mG-P_m|\to-5\log\phiG,
\end{equation}
and Zudilin's denominator bounds, in the convenient form
\begin{equation}\label{eq:Zden}
 2^{e_m}Q_m\in\Z,
 \qquad
 2^{e_m}D_{2m-1}^{\,2}P_m\in\Z,
 \qquad
 e_m=4m+O(\log m).
\end{equation}

The second row is modular.  Put
\[
 t(\tau)=
 \frac{\eta(\tau)^4\eta(4\tau)^2\eta(8\tau)^4}{\eta(2\tau)^{10}},
 \qquad
 F(\tau)=
 \frac{\eta(2\tau)^{10}}{\eta(\tau)^4\eta(4\tau)^4}.
\]
Then $F=\sum_{n\ge0}A_nt^n$, where
\begin{equation}\label{eq:Erec}
 (n+1)^2u_{n+1}
 =\bigl(12n(n+1)+4\bigr)u_n-32n^2u_{n-1},
\end{equation}
with $A_0=1,A_1=4$.  The canonical modular source is
\[
 \Phi_E=F\,q\frac{dt}{dq}=S-8S(2\tau),
\]
where $S=E_{3,\chi_{-4},1}$ in Fourier normalization.  Hence
\[
 L(\Phi_E,s)=(1-8\,2^{-s})\beta(s)\zeta(s-2),
 \qquad
 L(\Phi_E,2)=\frac{G}{2}.
\]
This explains the choice of the modular form: weight $3$ is the natural
weight for the value $L(2,\chi_{-4})$, and the oldform factor $1-8V_2$
kills the unwanted endpoint at $s=3$ while retaining the Catalan period.

Let $B(t)=\sum B_nt^n$ be the normalized companion $B_0=0,B_1=1$.  The
modular rectification is
\begin{equation}\label{eq:companion}
 B(t(q))=F(q)D_q^{-2}\Phi_E(q),
 \qquad D_q=q\frac d{dq}.
\end{equation}
Two integrations immediately give
\begin{equation}\label{eq:Eden}
 A_n\in\Z,
 \qquad D_n^2B_n\in\Z.
\end{equation}
Indeed the coefficient of $t^n$ only sees Fourier indices $m\le n$, and
$D_q^{-2}$ contributes only $m^{-2}$.  The characteristic roots of
\eqref{eq:Erec} are $8$ and $4$; therefore
\begin{equation}\label{eq:Earch}
 \frac1n\log|A_n|\to\log8,
 \qquad
 \frac1n\log\left|\frac G2A_n-B_n\right|\to\log4.
\end{equation}

\section{The hidden $2$-adic bridge}

The crucial fact is that the two recurrences represent the same $2$-adic
period.  We isolate the exact statement needed below.

\begin{lemma}[common $2$-adic Catalan period]\label{lem:p2}
There exists $G_2\in\Qtwo$ such that
\begin{align}
 \vTwo\!\left(G_2-\frac{P_m}{Q_m}\right)
 &=8m-1-4s_2(m),\label{eq:Ztail}\\
 \vTwo\!\left(G_2-\frac{2B_n}{A_n}\right)
 &=5n-1-4s_2(n),\label{eq:Etail}
\end{align}
where $s_2(r)$ is the number of $1$'s in the binary expansion of $r$.
Consequently
\begin{equation}\label{eq:diag}
 \vTwo(P_mA_{2m}-2B_{2m}Q_m)=4m-1.
\end{equation}
\end{lemma}

\begin{proof}[Idea of proof]
The exact Wronskians are
\[
 P_mQ_{m-1}-Q_mP_{m-1}
 =(-1)^{m-1}\frac{20m^2-8m+1}{8m^2(2m-1)^2}
\]
and
\[
 A_nB_{n-1}-A_{n-1}B_n=-\frac{32^{n-1}}{n^2}.
\]
Together with
\[
 \vTwo(Q_m)=-4m+2s_2(m),
 \qquad
 \vTwo(A_n)=2s_2(n),
\]
they give the exact tail valuations \eqref{eq:Ztail} and \eqref{eq:Etail}
once the two limits are identified.

That identification is the meeting point of two classical constructions.
Rivoal's moving Pad\'e specialization of Beukers's continued fraction is
exactly Zudilin's sequence.  If one defines the Pad\'e target with the
sign-safe normalization
\[
 \mathcal T(x)=2R(x)-R(x/2),
\]
where $R$ is Beukers's $p$-adic Hurwitz series, then
\[
 \mathcal T(x+1)+\mathcal T(x)=\frac{2}{x^2},
 \qquad
 \mathcal T_2(1/2)=8\zeta_2(2).
\]
At the moving point $x_m=\tfrac12-m$, Rivoal's exact Pad\'e identity and the
shift equation show that $P_m/Q_m$ converges to the single constant
$\zeta_2(2)$.  Beukers's Pad\'e estimate is uniform for the odd numerator
$1-2m$, so the moving-point passage is legitimate.

On the modular side, \eqref{eq:companion} is the level-$8$ form of Calegari's
overconvergent $2$-adic Catalan construction: after the standard
Teichm\"uller normalization its constant term is the same $\zeta_2(2)$.
Thus both limits equal $G_2:=\zeta_2(2)$.  Finally, at $n=2m$ the two tail
valuations are respectively
\[
 8m-1-4s_2(m),\qquad10m-1-4s_2(m),
\]
so the first is strictly smaller.  Ultrametricity therefore gives equality
for the difference of the two ratios; multiplying by $Q_mA_{2m}$ yields
\eqref{eq:diag}.
\end{proof}

The diagonal identity is more useful than it first appears, because it forces
the two $2$-adic limits to be equal; hence the exact tails give an
off-diagonal inequality for arbitrary indices.  If
\[
 \Delta_{i,k}=P_iA_k-2Q_iB_k,
\]
then
\begin{equation}\label{eq:roof}
\begin{split}
 \vTwo(\Delta_{i,k})\ge\min\{&
 4i-1-2s_2(i)+2s_2(k),\\
 &5k-4i-1+2s_2(i)-2s_2(k)\}.
\end{split}
\end{equation}
The two affine slopes balance when $8i=5k$.  This is the arithmetic reason
for the apparently strange sampling ratio
\[
 \boxed{i=5n,\qquad k=8n.}
\]
Since $s_2(8n)=s_2(n)$, \eqref{eq:roof} gives
\begin{equation}\label{eq:58delta}
 \vTwo(\Delta_{5n,8n})
 \ge20n-1-2|s_2(5n)-s_2(n)|
 =20n-O(\log n).
\end{equation}

\section{The correlated congruence lattice}

Put
\[
 S_n=D_{10n}^2.
\]
Since $10n$ is never a prime power, $D_{10n}=D_{10n-1}$.  Define the two
integer rows
\[
 (X_{1,n},Y_{1,n})
 =S_n(A_{8n},2B_{8n}),
\]
\[
 (X_{2,n},Y_{2,n})
 =2^{e_{5n}}S_n(Q_{5n},P_{5n}).
\]
Write $X_{r,n}=S_na_{r,n}$.  Their mixed minor after removing the common
$S_n$ from the first coordinates is
\[
 h_n=a_{1,n}Y_{2,n}-a_{2,n}Y_{1,n}
 =2^{e_{5n}}S_n\Delta_{5n,8n}.
\]
Since $e_{5n}=20n+O(\log n)$ and \eqref{eq:58delta} holds, there is an
explicit power $T_n$ with
\begin{equation}\label{eq:T}
 T_n\mid h_n,
 \qquad
 \frac1n\log T_n\to40\log2.
\end{equation}
For example one may take
\[
 T_n=2^{40n-3-2\lfloor\log_2(5n)\rfloor}.
\]

The point is that this determinant divisibility is \emph{correlated}: it can
be used in two congruences at the cost of one copy of $T_n$, rather than
$T_n^2$.  Set $M_n=S_nT_n$ and consider
\[
 K_n=\left\{(c_1,c_2)\in\Z^2:
 \begin{array}{l}
 a_{1,n}c_1+a_{2,n}c_2\equiv0\pmod{T_n},\\
 Y_{1,n}c_1+Y_{2,n}c_2\equiv0\pmod{S_nT_n}
 \end{array}
 \right\}.
\]
Then
\[
 q(c)=\frac{c_1X_{1,n}+c_2X_{2,n}}{M_n},
 \qquad
 p(c)=\frac{c_1Y_{1,n}+c_2Y_{2,n}}{M_n}
\]
are integers.

A two-line Smith-normal-form computation gives the useful index formula
\[
 [\Z^2:K_n]
 =\frac{S_nT_n^2}
 {\gcd(h_n,T_nY_{1,n},T_nY_{2,n},S_nT_na_{1,n},S_nT_na_{2,n},S_nT_n^2)}.
\]
Since $T_n\mid h_n$, we obtain
\begin{equation}\label{eq:index}
 [\Z^2:K_n]\le S_nT_n=M_n.
\end{equation}
By replacing one basis vector by a suitable positive integer multiple we may,
without losing integrality, pass to a sublattice $L_n\subseteq K_n$ with
\begin{equation}\label{eq:covol}
 M_n\le\operatorname{covol}(L_n)<2M_n.
\end{equation}
Thus the division modulus and the coefficient-lattice covolume have the same
exponential rate
\begin{equation}\label{eq:sigma}
 \sigma=20+40\log2,
\end{equation}
because $\log D_N=N+o(N)$.

\section{Balancing the two rows}

Let
\[
 \lambda_{r,n}=X_{r,n}G-Y_{r,n}.
\]
From \eqref{eq:Zarch}--\eqref{eq:Earch} and $\log S_n=20n+o(n)$ we get
\begin{align*}
 A_1&=20+24\log2,& E_1&=20+16\log2,\\
 A_2&=20+20\log2+25\log\phiG,
 &E_2&=20+20\log2-25\log\phiG,
\end{align*}
where $A_r$ is the exponential height of $X_{r,n}$ and $E_r$ that of
$\lambda_{r,n}$.  Moreover
\[
 (A_2+E_1)-(A_1+E_2)=50\log\phiG-8\log2>0,
\]
so the mixed determinant has a unique dominant real term.

We use the following standard rank-two form of the two-successive-minima
argument.  If a coefficient lattice has covolume $e^{\sigma n+o(n)}$, the
outputs are divided by $e^{\sigma n+o(n)}$, and the above crossed gap holds,
then choose the coefficient anisotropy
\[
 x=\frac{\sigma+E_2-E_1}{2}.
\]
Minkowski's second theorem applied to the rectangle
$|c_1|\le e^{xn}$, $|c_2|\le e^{(\sigma-x)n}$ yields an integer output with
\begin{equation}\label{eq:HFgen}
 \log|q_n|\ge Hn-o(n),
 \qquad
 \log|q_nG-p_n|\le Fn+o(n),
\end{equation}
where
\[
 H=A_2-x,
 \qquad
 F=x+E_1-\sigma.
\]
The use of two successive minima, rather than only Minkowski's first theorem,
is what prevents an exceptionally short rationality-kernel vector from
spoiling the lower bound on $|q_n|$.

Substituting \eqref{eq:sigma} gives
\[
 x=10+22\log2-\frac{25}{2}\log\phiG,
\]
and therefore
\begin{equation}\label{eq:HF}
 \boxed{H=10-2\log2+\frac{75}{2}\log\phiG},
 \qquad
 \boxed{F=10-2\log2-\frac{25}{2}\log\phiG}.
\end{equation}
Both are positive, and $H-F=50\log\phiG$.

\begin{theorem}\label{thm:main}
There are integer pairs $(q_n,p_n)$ with $q_n\ne0$, $q_nG-p_n\ne0$ and
$|q_n|\to\infty$ such that
\[
 \liminf_{n\to\infty}
 \frac{-\log|G-p_n/q_n|}{\log|q_n|}
 \ge
 \frac{50\log\phiG}
 {10-2\log2+\frac{75}{2}\log\phiG}
 =0.9025266028569714694\ldots .
\]
\end{theorem}

\begin{proof}
Apply \eqref{eq:HFgen} and use
\[
 -\log|G-p_n/q_n|
 =\log|q_n|-\log|q_nG-p_n|.
\]
Hence the liminf is at least $1-F/H=(H-F)/H$, and \eqref{eq:HF} gives the
stated constant.
\end{proof}

\section{Why $5:8$ is the scalar optimum}

There is a useful conceptual coda.  Sample Zudilin at $i=an$ and the modular
row at $k=cn$.  The common LCM-square contributes $2\max(2a,c)$ to the
exponential logarithmic cost, while the $2$-adic roof \eqref{eq:roof}
contributes
\[
 \min(8a,5c)\log2.
\]
The same balancing calculation gives
\[
 H(a,c)=\max(2a,c)+(2a+c)\log2
 -\frac12\min(8a,5c)\log2
 +\frac{15}{2}a\log\phiG,
\]
and
\[
 \delta(a,c)=\frac{10a\log\phiG}{H(a,c)}.
\]
By homogeneity put $r=c/a$.  Then $H(r)$ is strictly decreasing for
$0<r<8/5$ and strictly increasing for $r>8/5$ (with the elementary change of
formula at $r=2$).  Thus
\[
 \boxed{r=\frac85}
\]
is the unique optimum within this specific scalar architecture.  In other
words, the same equality $8i=5k$ that balances the two $2$-adic convergence
rates also optimizes the final Diophantine exponent.

\begin{remark}
The theorem is deliberately a statement about the \emph{worthiness} of an
explicit family of rational approximations.  Since the bound is $<1$, it does
not imply that $G$ is irrational.  The scalar optimum above is likewise only
an optimum inside the present two-realization, common-LCM-square,
$2$-adic-cross, two-minimum construction; it is not a universal obstruction
to other Catalan constructions.
\end{remark}

\begin{thebibliography}{9}

\bibitem{Zudilin2002}
W.~Zudilin,
\emph{An Ap\'ery-like difference equation for Catalan's constant},
arXiv:math/0201024 (2002).

\bibitem{ZudilinRemarks}
W.~Zudilin,
\emph{A few remarks on linear forms involving Catalan's constant},
arXiv:math/0210423 (2002).

\bibitem{Rivoal2006}
T.~Rivoal,
\emph{Nombres d'Euler, approximants de Pad\'e et constante de Catalan},
Ramanujan J. \textbf{11} (2006), 199--214.

\bibitem{Calegari2004}
F.~Calegari,
\emph{Irrationality of certain $p$-adic periods for small $p$},
arXiv:math/0408214.

\bibitem{Beukers2006}
F.~Beukers,
\emph{Irrationality of some $p$-adic $L$-values},
arXiv:math/0603277.

\end{thebibliography}

\end{document}
